%==============================================================================
%  CAPÍTULO VII
%==============================================================================
\chapter{Procedimiento concursal de liquidación ordinario}
\label{cap:liquidacion-ordinaria}

\fichaCapitulo{Realización colectiva de activos empresariales.}{Desasimiento, juicio de
oposición, causales de liquidación forzosa y régimen recursivo restrictivo.}

%==============================================================================
\bloqueTeorico
%==============================================================================

\subsection{La liquidación como \lat{ultima ratio} del derecho de insolvencia}
\pendiente{Desarrollar.}

\subsection{Desasimiento y administración privativa del liquidador}
\pendiente{Desarrollar la tensión entre realización colectiva y ejecuciones individuales
de acreedores con garantía real.}

%==============================================================================
\bloqueNormativo
%==============================================================================

\subsection{La liquidación voluntaria ordinaria}

Requisitos de entrada del \art{115}: carpetas tributarias completas e inventario
certificado de activos.

\subsection{La liquidación forzosa ordinaria}

Análisis de las causales del \art{117}, con especial atención a la cesación de pagos de
obligaciones derivadas del giro comercial.\verificar{Verificar número y contenido de las
causales en el texto vigente.}

\subsection{Desasimiento e incautación}

Régimen del \art{164}.

\subsection{Cuenta final de administración y sobreseimiento definitivo}
\pendiente{Desarrollar.}

\subsection{El régimen recursivo restrictivo}

El \art{4} y la apelación como excepción, concedida en el solo efecto devolutivo.

%==============================================================================
\bloquePractico
%==============================================================================

\subsection{Defensa del deudor en el juicio de oposición}

\begin{enumerate}
  \item \textbf{Notificación de la demanda de liquidación forzosa.} Emplazamiento legal
  de la empresa deudora.

  \item \textbf{Preparación de la audiencia inicial.} El deudor dispone de plazo fatal
  para consignar fondos que enerven la causal de cobro, o bien para preparar la prueba de
  su solvencia.\verificar{Confirmar el plazo aplicable y su cómputo.}

  \item \textbf{Audiencia inicial.} Comparecencia obligatoria del deudor y del acreedor
  demandante. El deudor puede oponerse formulando oralmente sus excepciones.

  \item \textbf{Audiencia de prueba.} Si existen puntos sustanciales controvertidos, se
  abre término probatorio y se rinde prueba testimonial y documental.\verificar{Duración
  del término probatorio.}

  \item \textbf{Audiencia de fallo.} Dictación de la sentencia definitiva. Acogida la
  demanda, se ordena la incautación de activos y el deudor pierde la administración.
\end{enumerate}

\begin{alerta}[Concentración procesal]
La estructura de audiencias sucesivas y plazos breves deja escaso margen de reacción.
La defensa debe estar preparada antes de la audiencia inicial: la prueba de solvencia no
se improvisa.
\end{alerta}

%==============================================================================
\bloqueJurisprudencial
%==============================================================================

\subsection{Constitucionalidad de la restricción del recurso de apelación}

Sistematización de las sentencias del \tconst{} dictadas en requerimientos de
inaplicabilidad por inconstitucionalidad ---entre ellas los \rol{12.539-2021-INA} y
\rol{12.527-2021-INA}---, y evolución del criterio del órgano de control constitucional
sobre si la denegación de la doble instancia vulnera la igual protección de la ley en el
ejercicio de los derechos y el debido proceso.\verificar{Confirmar roles, fechas y
resultado de cada requerimiento.}
